\documentclass[aspectratio=169,11pt,t]{beamer}
\input{../../../../theme/ksa-cs-v2.tex}
% 원본 덱 (ksa-theme): slides/v2/ml1/week15/kl_divergence/ -- 그림은 거기 figs/ 를 그대로 쓴다
\setbeamerfont{author}{size=\normalsize}
\setbeamerfont{institute}{size=\small}
\graphicspath{{../../../../v2/ml1/week15/kl_divergence/figs/}{../../../../assets/}}
\makeatletter
\newcommand{\markpage}[1]{\protected@write\@auxout{}{\string\newlabel{#1}{{\the\c@framenumber}{}}}\hypertarget{#1}{}}
\newcommand{\pref}[1]{\hyperlink{#1}{\textcolor{ksagolddk}{(p.\ref{#1})}}}
\makeatother

\renewcommand{\ksacredit}{Machine Learning and Deep Learning \,\textperiodcentered\, KL 발산}
\title{KL 발산 (Kullback–Leibler Divergence)}
\subtitle{믿은 확률과 실제 확률 사이의 거리 아닌 거리}
\author{Machine Learning and Deep Learning}
\institute{한국과학영재학교 (KSA)}
\date{}
\begin{document}
{\setbeamertemplate{footline}{}\begin{frame}[c,noframenumbering]\titlepage\end{frame}}


\section{질문}
% --- 오늘의 질문
\begin{frame}{오늘의 질문: 공지 확률이 틀렸다면 얼마나 손해일까?}
\vfill
\centering
\renewcommand{\arraystretch}{1.5}
{\Large\begin{tabular}{lcc}
 & 당첨 & 꽝\\\hline
공지 $Q$ & 20\% & 80\%\\
실제 $P$ & \kq{4\%} & 96\%
\end{tabular}}\vskip0.8em
{\large 당첨까지 평균 \kb{5}번 (공지) vs \kq{25}번 (실제): 돈이 5배}\\[0.4em]
{\large 두 확률 분포가 ``얼마나 다른지'' 를 숫자 하나로 잴 수 있을까?}\\[0.2em]
{\small 게임 뽑기는 설명용 예시 (실제 게임 아님)}
\vfill
\note{어떤 게임 뽑기의 공지 확률이 당첨 20퍼센트인데, 실제로는 4퍼센트였다고 해 봅시다. 설명을 위한 예시입니다. 공지대로라면 당첨까지 평균 5번이면 되지만, 실제로는 평균 25번, 돈이 5배 듭니다. 이렇게 믿은 확률과 실제 확률이 얼마나 다른지를 숫자 하나로 잴 수 있을까요? 그 숫자가 오늘 배울 KL 발산입니다. 분류 학습의 손실, VAE, t-SNE에서 이미 만났던 그 KL입니다.}
\end{frame}

\section{놀라움}
% --- 흉내
\begin{frame}{흉내 내 보기: 10만 명이 뽑으면}
\vfill
\centering\includegraphics[width=0.66\linewidth]{kl_gacha.png}\par\vskip0.2em
{\large 실제 4\%: 평균 \kq{24.9}번, 중앙값 17번 \quad 50번 넘게 \kq{13\%}, 100번 넘게 1.7\%}\\[0.2em]
{\small 공지 20\% 로 흉내 내면 평균 5.0번 \quad 평균 횟수 = 1 / 당첨 확률}
\vfill
\note{컴퓨터로 10만 명이 실제 확률 4퍼센트로 당첨될 때까지 뽑게 해 봤습니다. 평균 24.9번, 중앙값 17번이고, 50번 넘게 뽑은 사람이 13퍼센트, 100번 넘게 뽑은 사람도 1.7퍼센트입니다. 공지 확률 20퍼센트로 흉내 내면 평균 5.0번입니다. 당첨까지 평균 횟수는 1 나누기 당첨 확률이기 때문입니다. 이제 이 차이를 정보의 언어로 재 봅시다.}
\end{frame}

% --- 놀라움
\begin{frame}{놀라움 (정보량): $I = -\log_2 p$}
\vfill
\begin{columns}[c]
\column{0.5\linewidth}\centering
\includegraphics[width=\linewidth]{kl_surprise.png}
\column{0.48\linewidth}
{\large 드문 일일수록 놀랍다}\\[0.4em]
\renewcommand{\arraystretch}{1.3}
{\large\begin{tabular}{lcc}
 & 실제로 재면 & 공지로 재면\\\hline
당첨 & \kb{4.64} & 2.32\\
꽝 & 0.06 & 0.32
\end{tabular}}\\[0.6em]
{\large 단위: 비트 ($\log_2$) \quad 확률 1/2 = 1비트}\\
{\small 동전 한 번 결과를 적는 데 필요한 0 · 1 의 개수}
\end{columns}
\vfill
\note{먼저 놀라움, 정보량입니다. 마이너스 로그 2 p로 잽니다. 확률이 작을수록, 드문 일일수록 놀라움이 큽니다. 실제 확률 4퍼센트로 재면 당첨은 4.64비트, 꽝은 0.06비트입니다. 공지 확률 20퍼센트로 재면 당첨은 2.32비트, 꽝은 0.32비트입니다. 단위는 비트인데, 확률 2분의 1인 일, 동전 한 번의 결과가 1비트입니다. 그 결과를 0과 1로 적는 데 필요한 자리 수라고 생각하면 됩니다.}
\end{frame}

% --- 엔트로피 · 교차 엔트로피
\begin{frame}{엔트로피 · 교차 엔트로피}
\markpage{ce}
\vfill
\centering
\renewcommand{\arraystretch}{1.8}
{\large\begin{tabular}{lll}
 & 놀라움은 & 평균은 실제 비율로\\\hline
엔트로피 $H(P)$ & 실제 확률로 & $4.64 \times 0.04 + 0.06 \times 0.96 = \kb{0.24}$\\
교차 엔트로피 $H(P, Q)$ & \textbf{공지} 확률로 & $2.32 \times 0.04 + 0.32 \times 0.96 = \kq{0.40}$
\end{tabular}}\vskip0.6em
{\Large\[ H(P) = -\sum_x P(x) \log_2 P(x), \qquad H(P, Q) = -\sum_x P(x) \log_2 Q(x) \]}\vskip0.2em
{\large 뜻: 결과를 하나 적을 때 드는 평균 비트 \quad 틀린 확률로 부호를 만들면 더 든다}
\vfill
\note{엔트로피는 실제 확률로 잰 놀라움을 실제 비율로 평균 낸 값입니다. 4.64 곱하기 0.04 더하기 0.06 곱하기 0.96, 0.24비트입니다. 교차 엔트로피는 놀라움은 공지 확률로 재고, 평균은 실제 비율로 냅니다. 2.32 곱하기 0.04 더하기 0.32 곱하기 0.96, 0.40비트입니다. 뜻으로 보면, 뽑기 결과를 하나 적을 때 드는 평균 비트 수입니다. 실제 확률에 맞는 부호를 쓰면 0.24비트면 되지만, 틀린 공지 확률로 부호를 만들면 0.40비트가 듭니다.}
\end{frame}

\section{KL}
% --- KL
\begin{frame}{KL 발산 = 교차 엔트로피 $-$ 엔트로피}
\markpage{kl}
\vfill
\centering\includegraphics[width=0.46\linewidth]{kl_bar.png}\par\vskip0.2em
{\LARGE\[ \text{KL}(P \,\|\, Q) = \sum_x P(x) \log_2 \frac{P(x)}{Q(x)} = 0.40 - 0.24 = \kb{0.16} \text{ 비트} \]}\vskip0.1em
{\large 틀린 확률을 믿어서 \textbf{한 번에 더 드는} 비트 \quad 뽑기 1,000번이면 160비트 낭비}\\[0.2em]
{\large 항상 $\ge 0$, \quad $P = Q$ 일 때만 0 \quad (Kullback \& Leibler, 1951)}
\vfill
\note{교차 엔트로피에서 엔트로피를 빼면 KL 발산입니다. 0.40 빼기 0.24, 0.16비트입니다. 식으로는 P 곱하기 로그 P 나누기 Q의 합입니다. 뜻은 틀린 확률을 믿어서 결과 하나를 적을 때마다 더 드는 비트입니다. 뽑기 1000번을 기록하면 160비트를 낭비하는 셈이죠. KL은 항상 0 이상이고, 두 분포가 같을 때만 0입니다. 1951년 쿨백과 라이블러가 정의했습니다.}
\end{frame}

% --- 비대칭
\begin{frame}{거리처럼 보이지만 대칭이 아니다}
\markpage{asym}
\vfill
\centering
{\large $\text{KL}(P \| Q) = \kb{0.16}$ \qquad $\text{KL}(Q \| P) = 0.2 \log_2 \frac{0.2}{0.04} + 0.8 \log_2 \frac{0.8}{0.96} = \kq{0.25}$}\vskip0.4em
\includegraphics[width=0.84\linewidth]{kl_modes.png}\par
{\large 봉우리 둘인 $P$ 를 정규분포 $Q$ 로: \textcolor{ksablue}{$\text{KL}(P\|Q)$ 는 다 덮고}, \textcolor{ksared}{$\text{KL}(Q\|P)$ 는 하나만}}
\vfill
\note{KL은 거리처럼 보이지만 대칭이 아닙니다. P와 Q를 바꾸면 0.25비트로 달라집니다. 이 차이는 모델을 맞출 때 중요합니다. 봉우리가 둘인 분포 P를 정규분포 하나 Q로 맞춰 봤습니다. KL P Q를 줄이면 P가 있는 곳은 Q가 꼭 덮어야 하니, 평균 0, 표준편차 2.06으로 두 봉우리를 다 덮습니다. KL Q P를 줄이면 Q가 있는 곳에 P가 있어야 하니, 봉우리 하나를 골라 거기에만 붙습니다. 첫 번째를 다 덮기, 두 번째를 하나 고르기라고 합니다. GAN이 한 모양만 만드는 모드 붕괴와도 이어지는 이야기입니다.}
\end{frame}

\section{쓰임}
% --- ML 속 KL
\begin{frame}{머신러닝 속 KL}
\vfill
\centering
\renewcommand{\arraystretch}{1.5}
{\large\begin{tabular}{ll}
분류 손실 (2주차) & 정답이 원-핫이면 교차 엔트로피 = KL\\
VAE (15주차) & 구름을 $N(0, 1)$ 쪽으로: $\text{KL}(N(\mu, \sigma^2) \| N(0,1))$\\
t-SNE (14주차) & 고차원 이웃 $P$ 와 지도 이웃 $Q$: $\text{KL}(P \| Q)$\\
지식 증류 & 작은 모델이 큰 모델의 확률을 따라 하게: KL\\
ChatGPT 의 RLHF & 원래 모델에서 멀어지지 않게 KL 벌점
\end{tabular}}\\[0.6em]
{\large 분류: 고양이 정답에 70\% $\to$ 0.51 비트, 95\% $\to$ \kb{0.07} 비트}
\vfill
\note{KL은 머신러닝 곳곳에 있습니다. 2주차에 배운 분류의 교차 엔트로피 손실은, 정답이 원-핫이면 엔트로피가 0이라 KL과 같습니다. 고양이 사진에 70퍼센트를 주면 0.51비트, 95퍼센트를 주면 0.07비트입니다. VAE는 구름을 N 0 1 쪽으로 당기는 데, t-SNE는 고차원과 지도의 이웃 확률을 맞추는 데 KL을 썼습니다. 큰 모델의 확률을 작은 모델이 따라 하게 하는 지식 증류도 KL이고, ChatGPT를 사람 선호에 맞추는 RLHF는 원래 모델에서 너무 멀어지지 않도록 KL 벌점을 줘서, 점수만 높이는 꼼수, 보상 해킹을 막습니다.}
\end{frame}

\section{연습}
% --- 문제풀이
\begin{frame}{문제풀이}
\vfill
\begin{enumerate}
\item 공정한 동전 $P = (0.5, 0.5)$ 을 $Q = (0.9, 0.1)$ 로 믿으면 $\text{KL}(P\|Q)$ 는? ($\log_2 \frac{0.5}{0.9} \approx -0.848$, $\log_2 5 \approx 2.32$)\vskip0.8em
\item 정답 (1, 0) 에 모델이 (0.5, 0.5) 를 주면 손실 (비트) 은?\vskip0.8em
\item 당첨 확률이 2\% 라면 당첨까지 평균 몇 번? 놀라움은 몇 비트? ($\log_2 50 \approx 5.64$)\vskip0.8em
\item 데이터 $P$ 를 흉내 내는 생성 모델이 $\text{KL}(Q\|P)$ 만 줄이면 어떤 문제가 생길까?
\end{enumerate}
\vfill
\note{1번 답. 0.5 곱하기 마이너스 0.848 더하기 0.5 곱하기 2.32, 약 0.74비트입니다. 2번 답. 마이너스 로그 2 0.5, 1비트입니다. 3번 답. 1 나누기 0.02, 50번이고, 놀라움은 5.64비트입니다. 4번 답. 봉우리 하나만 골라도 KL Q P가 작으니, 데이터의 여러 모양 중 일부만 만들어 내는 모드 붕괴가 생길 수 있습니다.}
\end{frame}

% --- 실습
\begin{frame}{실습}
\vfill
\centering\begin{tabular}{rl}
\large 1 & 놀라움 · 엔트로피 · 교차 엔트로피 · KL 함수 $\to$ 0.24 · 0.40 · 0.16\\[0.5em]
\large 2 & 뽑기 흉내: 당첨까지 횟수, 기록 1,000번의 비트 수\\[0.5em]
\large 3 & 봉우리 둘을 정규분포로 맞추기: 앞 방향 vs 뒤 방향 KL
\end{tabular}\vskip1em
\href{https://colab.research.google.com/github/SuminHan/book-ml/blob/main/slides/v2/ml1/week15/kl_divergence/lab/kl_divergence.ipynb}{\textcolor{ksagolddk}{▶ Colab에서 실습 노트북 열기}}
\vfill
\note{실습은 세 부분입니다. 첫째, 놀라움, 엔트로피, 교차 엔트로피, KL을 함수로 만들어 0.24, 0.40, 0.16을 확인합니다. 둘째, 뽑기를 흉내 내 당첨까지 횟수를 세고, 결과 1000번을 기록하는 비트 수를 비교합니다. 셋째, 봉우리 둘인 분포를 정규분포 하나로 맞추며 앞 방향과 뒤 방향 KL을 비교합니다.}
\end{frame}

\section{정리}
% --- 정리
\begin{frame}{정리}
\vfill
\centering\large\begin{tabular}{l}
놀라움 $-\log_2 p$ \quad 엔트로피 = 평균 놀라움 (뽑기 0.24 비트)\\[0.5em]
교차 엔트로피 = 틀린 확률로 잰 평균 놀라움 (0.40) \quad \textbf{KL} = 그 차이 \kb{0.16}\\[0.5em]
KL $\ge 0$, 같을 때만 0 \quad 대칭 아님 (0.16 vs 0.25): 다 덮기 vs 하나 고르기\\[0.5em]
분류 손실 · VAE · t-SNE · 지식 증류 · RLHF 의 공통 언어
\end{tabular}
\vfill
\note{정리합니다. 놀라움은 마이너스 로그 확률이고, 엔트로피는 평균 놀라움입니다. 뽑기에서는 0.24비트였습니다. 교차 엔트로피는 틀린 확률로 잰 평균 놀라움 0.40비트이고, 그 차이 0.16비트가 KL 발산입니다. KL은 항상 0 이상이고 같을 때만 0이며, 대칭이 아니라서 다 덮기와 하나 고르기의 차이가 생깁니다. 분류 손실, VAE, t-SNE, 지식 증류, RLHF가 모두 이 언어로 쓰여 있습니다.}
\end{frame}
\end{document}
